% concatenated from symmetric_polar_degree_paper_updated.tex
\documentclass[11pt]{article}

% ---------------------------------------------------------------
% Formatted to match the structure and conventions of
% computational complexity (Springer/Birkhauser).  The file uses
% the standard article class for portability in this environment;
% the body is intended to transfer directly into a Springer Nature
% journal template for submission.
% ---------------------------------------------------------------

\usepackage[margin=1in]{geometry}
\usepackage{amsmath,amssymb,amsthm}
\usepackage{mathtools}
\usepackage{bm}
\usepackage[hidelinks]{hyperref}
\usepackage{enumitem}
\usepackage{microtype}
\usepackage{upquote}

\theoremstyle{plain}
\newtheorem{theorem}{Theorem}
\newtheorem{lemma}{Lemma}
\newtheorem{proposition}{Proposition}
\newtheorem{corollary}{Corollary}
\theoremstyle{definition}
\newtheorem{definition}{Definition}
\newtheorem{remark}{Remark}

\DeclareMathOperator{\sdc}{sdc}
\DeclareMathOperator{\dc}{dc}
\DeclareMathOperator{\rank}{rank}
\DeclareMathOperator{\adj}{adj}
\DeclareMathOperator{\tr}{tr}
\DeclareMathOperator{\Det}{Det}
\DeclareMathOperator{\codim}{codim}
\DeclareMathOperator{\Spec}{Spec}

\newcommand{\PP}{\mathbb{P}}
\newcommand{\CC}{\mathbb{C}}
\newcommand{\OO}{\mathcal O}
\newcommand{\Mcal}{\mathcal M}
\newcommand{\pdeg}{\delta}
\newcommand{\veps}{\varepsilon}

\title{\bf A symmetric determinantal lower bound for diagonal power sums\\
via polar degree}

\author{%
  Karthik Sheshadri\thanks{Independent AI researcher and engineer, San Jose,
  California, USA. The author holds a PhD in computer science and has
  graduate-level training in mathematics, but is not a specialist in algebraic
  complexity or algebraic geometry; see the ``Status of this work'' note
  following the abstract. Correspondence:
  \texttt{karthiksheshadri217@gmail.com}. LinkedIn:
  \url{https://www.linkedin.com/in/karthik-sheshadri-0624ab150/}.}
  \\[4pt]
  \small Independent AI researcher and engineer, San Jose, California, USA
}

\date{}

\begin{document}
\maketitle

\begin{abstract}
The symmetric determinantal complexity $\sdc(f)$ of a polynomial $f$ is the
least $m$ such that $f=\Det(M)$ for an $m\times m$ symmetric matrix $M$ of
affine-linear forms.  We prove, over $\CC$, that
\[
  \sdc\!\left(\sum_{i=1}^n x_i^n\right)
  \ge \left(\frac{1}{2e}-o(1)\right)n^2 .
\]
The result is a symmetric companion to the author's non-symmetric polar-degree
preprint~\cite{SheshadriArxiv}.  The method parallels that work, but the proof
below is self-contained and redoes the load-bearing local incidence analysis in
the symmetric setting.  The general theorem is the following.  If
$X=V(f)\subset\PP^{N-1}$ is a smooth degree-$d$ hypersurface, $N\ge3$, and
$f=\Det(A_0+\sum_{i=1}^N x_iA_i)$ with all $A_i$ symmetric of size $m$, then
\[
  \pdeg_{\mathrm{top}}(X)=d(d-1)^{N-2}
  \le 2^{N-2}\binom{m}{N-1}.
\]
The proof uses the symmetric rank-one kernel incidence
$\Mcal(z,x)u=0$, where $\Mcal=zA_0+\sum_i x_iA_i$.  At a genuine polar point,
$\Mcal$ has rank $m-1$, and the symmetric local normal form
\[
  \Mcal=\begin{pmatrix}B&c\\ c^{\mathsf T}&s\end{pmatrix},\qquad
  \det B\in\OO^\times,
\]
eliminates the unique projective kernel line scheme-theoretically:
$u=(-B^{-1}c,1)$ and $\det\Mcal=(\det B)(s-c^{\mathsf T}B^{-1}c)$.  On this
local graph, $\adj(\Mcal)=(\det B)uu^{\mathsf T}$ along the determinant
hypersurface, so the lifted conormal forms $u^{\mathsf T}A_i u$ are a common
unit multiple of the ordinary partial derivatives $\partial_i f$.  Hence the
lifted polar equations cut the ordinary polar slice, up to units, and every
genuine lifted polar point is a zero-dimensional scheme-theoretic isolated
solution.  Multihomogeneous Bezout on
$\PP^N_{[z:x]}\times\PP^{m-1}_{[u]}$ then gives
\[
  [H^N U^{m-1}]\,H(H+U)^m(2U)^{N-2}
  =2^{N-2}\binom{m}{N-1}.
\]
For $F_n=\sum_i x_i^n$ this bounds $n(n-1)^{n-2}$ and yields the stated
constant $1/(2e)$.  More generally, for
$F_{N,d}=\sum_{i=1}^N x_i^d$ the same theorem gives
$\sdc(F_{N,d})\ge(1/(2e)-o_N(1))N(d-1)$ as $N\to\infty$, uniformly for
$d\ge2$.  We also give an explicit symmetric determinantal representation of
$F_{N,d}$ of size $2N(d+1)+1$, showing that the diagonal lower bounds are
non-vacuous and tight up to a constant factor.  The result is for exact
symmetric determinantal complexity in characteristic zero; it is not a
border-complexity statement and it is not a uniform positive-characteristic
theorem.
\end{abstract}

\noindent\textbf{Keywords:} symmetric determinantal complexity; polar degree;
Gauss map; dual variety; algebraic complexity; multihomogeneous Bezout;
Valiant's hypothesis.

\medskip
\noindent\textbf{Subject classification:} 68Q17; 14N05; 14Q20; 15A15.

\medskip
\noindent\textbf{Status of this work.}
Three points should be stated plainly.
\emph{(i) Provenance.}  This paper is a symmetric companion to the author's
non-symmetric polar-degree preprint~\cite{SheshadriArxiv}.  The method parallels
that preprint, but the proof below is self-contained: in particular, the
symmetric isolatedness step is reproved directly by a Schur-complement normal
form rather than imported from the non-symmetric argument.  The present statement
arose from a follow-up adversarial prompt asking whether the symmetric
specialization of the incidence argument is valid or whether the coincidence of
the left and right kernels introduces a new failure mode.  The prompt protocol is
described in Section~\ref{sec:methodology} and Appendix~\ref{app:prompt}.
\emph{(ii) What is proved.}  The main inequality, the isolatedness lemma, the
multidegree computation, and the $1/(2e)$ asymptotic extraction are proved in
full below.  The proof is not a heuristic specialization of the non-symmetric
case: the symmetric Schur-complement normal form is used directly to identify
the completed local incidence ring at a genuine polar point with the completed
local ring of the ordinary polar slice, up to units.
\emph{(iii) Scope.}  Everything is over $\CC$, or equivalently any algebraically
closed field of characteristic zero for the arguments used here.  The theorem is
false if stated uniformly over all fields of characteristic different from two;
see Remark~\ref{rem:positivechar}.

\section{Introduction}
\label{sec:intro}

\subsection{Symmetric determinantal complexity}

Let $f\in\CC[x_1,\ldots,x_N]$.  Its \emph{symmetric determinantal complexity},
denoted $\sdc(f)$, is the least integer $m$ for which there exist symmetric
matrices $A_0,A_1,\ldots,A_N\in\CC^{m\times m}$ such that
\[
  f(x)=\Det\!\left(A_0+\sum_{i=1}^N x_iA_i\right).
\]
This is a restriction of ordinary affine determinantal complexity $\dc(f)$, in
which the representing matrix need not be symmetric.  Determinantal complexity
is one of the central algebraic models in Valiant's program~\cite{Valiant79};
quadratic lower bounds for the permanent are known~\cite{MignonRessayre04}, but
superquadratic and superpolynomial lower bounds remain major open problems.

The ordinary determinant model has two kernel variables at a smooth point of
the determinant hypersurface: a left kernel line and a right kernel line.  The
symmetric model has only one.  At rank $m-1$ the left and right kernels coincide,
so the conormal direction of a symmetric determinant is controlled by quadratic
forms
\[
  u^{\mathsf T}A_i u
\]
in one projective kernel variable $[u]\in\PP^{m-1}$ rather than by bilinear
forms $u^{\mathsf T}A_i v$ in two projective kernel variables.  This reduction
is exactly what improves the Bezout constant from $1/(4e)$ in the ordinary
polar-degree count of the companion preprint~\cite{SheshadriArxiv} to
$1/(2e)$ here.

The improvement is model-specific.  Since $\sdc(f)\ge \dc(f)$, any lower bound
for $\dc(f)$ is automatically a lower bound for $\sdc(f)$, but the argument in
this paper is sharper because it uses the symmetric geometry directly.

\subsection{Main results}

The key geometric input is the top polar degree of a smooth hypersurface.  If
$X=V(f)\subset\PP^{N-1}$ is smooth of degree $d\ge2$, then the gradient map
\[
  \gamma_X:X\longrightarrow (\PP^{N-1})^\vee,
  \qquad
  [x]\longmapsto [\partial_1f(x):\cdots:\partial_N f(x)]
\]
is a morphism and $\gamma_X^*\OO(1)=\OO_X(d-1)$.  We define
$\pdeg_{\mathrm{top}}(X)$ as the length of the inverse image of a general
codimension-$(N-2)$ linear subspace of the dual projective space.  Equivalently,
for a smooth hypersurface,
\[
  \pdeg_{\mathrm{top}}(X)=\int_X c_1(\OO_X(d-1))^{N-2}=d(d-1)^{N-2}.
\]
This is also the classical dual-class degree in characteristic zero, but the
proof uses the polar intersection length, not any closedness property of dual
varieties.

\begin{theorem}[Symmetric polar-degree bound]
\label{thm:main}
Let $N\ge3$, and let $f\in\CC[x_1,\ldots,x_N]$ be homogeneous of degree
$d\ge2$ such that $X=V(f)\subset\PP^{N-1}$ is smooth.  Suppose
\[
  f(x)=\Det\!\left(A_0+\sum_{i=1}^N x_iA_i\right)
\]
for symmetric $m\times m$ complex matrices $A_i$.  Then
\[
  \pdeg_{\mathrm{top}}(X)=d(d-1)^{N-2}
  \le 2^{N-2}\binom{m}{N-1}.
\]
\end{theorem}


\begin{corollary}[Degree-$n$ diagonal power sums]
\label{cor:fermat}
For $F_n=x_1^n+\cdots+x_n^n$ over $\CC$,
\[
  \sdc(F_n)\ge \left(\frac{1}{2e}-o(1)\right)n^2.
\]
\end{corollary}

\begin{corollary}[General diagonal power sums]
\label{cor:diagonal-general}
Let
\[
  F_{N,d}=x_1^d+\cdots+x_N^d
\]
over $\CC$, with $N\ge3$ and $d\ge2$.  Then
\[
  d(d-1)^{N-2}
  \le 2^{N-2}\binom{\sdc(F_{N,d})}{N-1},
\]
and hence
\[
  \sdc(F_{N,d})
  \ge
  \left(
  \frac{d(d-1)^{N-2}(N-1)!}{2^{N-2}}
  \right)^{1/(N-1)}.
\]
In particular, uniformly for $d\ge2$ as $N\to\infty$,
\[
  \sdc(F_{N,d})
  \ge \left(\frac{1}{2e}-o_N(1)\right)N(d-1).
\]
\end{corollary}

The $n$-variable degree-$n$ case is stated separately because it is the most
compact asymptotic form and parallels the companion non-symmetric result.

\subsection{Why the symmetric isolatedness step is the whole proof}

The intended incidence count is visually simple.  Homogenize the representing
matrix:
\[
  \Mcal(z,x)=zA_0+\sum_{i=1}^N x_iA_i,
  \qquad \Det\Mcal=z^{m-d}f(x).
\]
A genuine polar point of $X$ is lifted to a point
$([z:x],[u])\in\PP^N\times\PP^{m-1}$ satisfying
\[
  \Mcal(z,x)u=0,
  \qquad
  q_j(u^{\mathsf T}A_1u,\ldots,u^{\mathsf T}A_Nu)=0,
  \qquad
  h(z,x)=0.
\]
The $m$ kernel equations have class $H+U$, the $N-2$ polar equations have class
$2U$, and the slice has class $H$.  Thus a square multihomogeneous Bezout count
would give
\[
  [H^N U^{m-1}]\,H(H+U)^m(2U)^{N-2}
  =2^{N-2}\binom{m}{N-1}.
\]

But Bezout bounds only isolated zero-dimensional contributions.  The
load-bearing issue is therefore not the coefficient extraction; it is whether a
genuine lifted polar point can lie on a positive-dimensional component of the
symmetric incidence scheme.  In particular, one must rule out the possibility
that the kernel equation acquires extra fiber dimension near a rank-drop point,
or that the quadratic forms $u^{\mathsf T}A_i u$ vanish in a way that creates a
spurious component through a genuine point.  Section~\ref{sec:local} proves
precisely this local isolatedness statement by a symmetric Schur-complement
normal form.  This is the point at which the symmetry constraint must be checked
rather than assumed.

\section{Preliminaries}
\label{sec:prelim}

We work over $\CC$.  Let
\[
  M(x)=A_0+\sum_{i=1}^N x_iA_i
\]
with all $A_i$ symmetric.  Its homogenization is
\[
  \Mcal(z,x)=zA_0+\sum_{i=1}^N x_iA_i.
\]
Since $\Det M(x)=f(x)$ is homogeneous of degree $d$, one has
\begin{equation}
\label{eq:homdet}
  \Det\Mcal(z,x)=z^{m-d}f(x).
\end{equation}
Here necessarily $m\ge d$ unless $f=0$, because the determinant of an affine
$m\times m$ matrix has degree at most $m$.

We use the identity
\begin{equation}
\label{eq:detderiv}
  \frac{\partial}{\partial x_i}\Det\Mcal
  =\tr\!\big(\adj(\Mcal)A_i\big).
\end{equation}
At a symmetric rank-$(m-1)$ matrix, the kernel is a line and the adjugate is a
nonzero symmetric rank-one matrix with image equal to that kernel line.

\begin{lemma}[Rank at genuine smooth points]
\label{lem:rank}
Let $[z_0:x_0]\in\PP^N$ satisfy $z_0\ne0$ and $f(x_0)=0$, with
$[x_0]\in X$ smooth.  Then
\[
  \rank\Mcal(z_0,x_0)=m-1.
\]
Consequently the projective kernel line $[u_0]\in\PP^{m-1}$ is unique.
\end{lemma}

\begin{proof}
Since $\Det\Mcal(z_0,x_0)=0$, the rank is at most $m-1$.  If the rank were at
most $m-2$, then all $(m-1)\times(m-1)$ cofactors would vanish, so
$\adj(\Mcal(z_0,x_0))=0$.  By \eqref{eq:detderiv}, every partial derivative of
$\Det\Mcal$ with respect to the $x_i$ would vanish at $[z_0:x_0]$.  But by
\eqref{eq:homdet}, and because $z_0\ne0$,
\[
  \frac{\partial}{\partial x_i}\Det\Mcal(z_0,x_0)
  =z_0^{m-d}\partial_i f(x_0).
\]
The point $[x_0]\in X$ is smooth, so not all $\partial_i f(x_0)$ vanish.  This
contradiction proves the rank is $m-1$.  A rank-$(m-1)$ matrix has one-dimensional
kernel.
\end{proof}

\begin{lemma}[Symmetric conormal identity]
\label{lem:conormal}
At a point as in Lemma~\ref{lem:rank}, let $0\ne u_0\in\ker\Mcal(z_0,x_0)$.  Then
there is a scalar $\alpha\ne0$ such that
\[
  \partial_i f(x_0)=\alpha\,u_0^{\mathsf T}A_i u_0
  \qquad (1\le i\le N).
\]
In particular,
\[
  [\partial_1f(x_0):\cdots:\partial_Nf(x_0)]
  =[u_0^{\mathsf T}A_1u_0:\cdots:u_0^{\mathsf T}A_Nu_0].
\]
\end{lemma}

\begin{proof}
At a symmetric rank-$(m-1)$ matrix, the adjugate is a nonzero scalar multiple of
$u_0u_0^{\mathsf T}$.  Thus
\[
  \adj(\Mcal(z_0,x_0))=\beta u_0u_0^{\mathsf T}
\]
for some $\beta\ne0$.  By \eqref{eq:detderiv},
\[
  \frac{\partial}{\partial x_i}\Det\Mcal(z_0,x_0)
  =\beta\,u_0^{\mathsf T}A_i u_0.
\]
Using \eqref{eq:homdet} and $z_0\ne0$ gives
$z_0^{m-d}\partial_i f(x_0)=\beta u_0^{\mathsf T}A_i u_0$.  Take
$\alpha=\beta z_0^{d-m}$.
\end{proof}

\section{The symmetric local normal form}
\label{sec:local}

This section proves the isolatedness statement needed for Bezout.  We state it
in the form used later.

Choose general linear forms $q_1,\ldots,q_{N-2}$ on the dual coordinates so that
\[
  P_q=\{[x]\in X:q_j(\partial_1f(x),\ldots,\partial_Nf(x))=0
  \text{ for }1\le j\le N-2\}
\]
is a reduced zero-dimensional scheme of length $\pdeg_{\mathrm{top}}(X)$.  This
is possible because the partial derivatives define a base-point-free linear
system $|\OO_X(d-1)|$ on the smooth variety $X$; iterated Bertini gives a
reduced complete intersection of degree $d(d-1)^{N-2}$.

Choose a linear form $\ell(x)$ nonzero on the finite set $P_q$ and set
\[
  h(z,x)=\ell(x)-z.
\]
Each $[\xi]\in P_q$ determines the affine-slice representative
$[z:x]=[\ell(\xi):\xi]$, which has $z\ne0$.

\begin{lemma}[Symmetric isolatedness]
\label{lem:isolated}
Let $p=([z_0:x_0],[u_0])$ be the lift of a point of $P_q$ to the symmetric
incidence scheme
\begin{equation}
\label{eq:incidence}
  \Mcal(z,x)u=0,
  \qquad
  q_j(u^{\mathsf T}A_1u,\ldots,u^{\mathsf T}A_Nu)=0\ (1\le j\le N-2),
  \qquad
  h(z,x)=0
\end{equation}
in $\PP^N_{[z:x]}\times\PP^{m-1}_{[u]}$.  Then $p$ is a zero-dimensional
scheme-theoretically isolated solution of \eqref{eq:incidence}.
\end{lemma}

\begin{proof}
By Lemma~\ref{lem:rank}, $\Mcal(z_0,x_0)$ has rank $m-1$.  Since it is
symmetric, its kernel is a one-dimensional radical.  After a constant congruence
change of basis $\Mcal\mapsto P^{\mathsf T}\Mcal P$, which gives an
isomorphic local incidence problem via $u=P u'$ and preserves the quadratic
forms since ${u'}^{\mathsf T}(P^{\mathsf T}A_iP)u'=u^{\mathsf T}A_i u$, we
may assume that the kernel line at $p$ is spanned by the last coordinate vector.
The determinant is multiplied only by the nonzero scalar $(\det P)^2$, which is
harmless in all local unit comparisons below.  Because the kernel is the radical
of the symmetric bilinear form, the restriction to any complementary
$(m-1)$-dimensional subspace is nondegenerate.  Thus, in a Zariski neighborhood
of $[z_0:x_0]$, write
\[
  \Mcal=\begin{pmatrix}B&c\\ c^{\mathsf T}&s\end{pmatrix},
  \qquad \det B\in\OO^\times.
\]
Work on the affine chart $u_m=1$ and write $u=(u',1)^{\mathsf T}$.  The kernel
equations are
\[
  Bu'+c=0,
  \qquad
  c^{\mathsf T}u'+s=0.
\]
Since $B$ is invertible in the local ring, the first $m-1$ equations are
equivalent to
\[
  u'=-B^{-1}c.
\]
With
\[
  g=s-c^{\mathsf T}B^{-1}c,
\]
the last equation becomes
\[
  c^{\mathsf T}u'+s
  =c^{\mathsf T}(u'+B^{-1}c)+g.
\]
Therefore the ideal generated by $\Mcal u$ is exactly
\[
  \big(u'+B^{-1}c,\ g\big)
\]
in the local ring.  Scheme-theoretically, the kernel incidence is the graph
\[
  u=\nu:=\begin{pmatrix}-B^{-1}c\\1\end{pmatrix}
\]
over the hypersurface $g=0$.  The determinant is
\begin{equation}
\label{eq:schurdet}
  \Det\Mcal=(\det B)g.
\end{equation}

On $g=0$, the matrix $\Mcal$ has kernel spanned by $\nu$ and rank $m-1$ because
$B$ remains invertible.  Its adjugate is therefore a scalar multiple of
$\nu\nu^{\mathsf T}$.  The lower-right cofactor equals $\det B$, while the
lower-right entry of $\nu\nu^{\mathsf T}$ is $1$, so
\begin{equation}
\label{eq:adjlocal}
  \adj(\Mcal)=(\det B)\nu\nu^{\mathsf T}
  \qquad\text{on }g=0.
\end{equation}
Hence, in the local quotient by the kernel equations,
\[
  \frac{\partial}{\partial x_i}\Det\Mcal
  =(\det B)\nu^{\mathsf T}A_i\nu.
\]
Using \eqref{eq:homdet}, we get
\begin{equation}
\label{eq:unitpolar}
  \nu^{\mathsf T}A_i\nu
  =\frac{z^{m-d}}{\det B}\,\partial_i f(x)
  \qquad\text{on the local kernel graph.}
\end{equation}
The factor $z^{m-d}/\det B$ is a unit near $p$, since $z_0\ne0$ and
$\det B$ is a unit.

It follows that, after eliminating the kernel variables, each lifted polar
equation
\[
  q_j(\nu^{\mathsf T}A_1\nu,\ldots,\nu^{\mathsf T}A_N\nu)=0
\]
is the ordinary polar equation
\[
  q_j(\partial_1f,\ldots,\partial_Nf)=0
\]
multiplied by the same local unit.  By \eqref{eq:schurdet} and
\eqref{eq:homdet}, the equation $g=0$ is also the equation $f=0$ up to a unit on
$z\ne0$.

Thus the completed local ring of the incidence scheme at $p$ is isomorphic to
\[
  \widehat\OO_{\PP^N,[z_0:x_0]}
  \Big/
  \big(f,\ q_1(\nabla f),\ldots,q_{N-2}(\nabla f),\ h\big),
\]
up to multiplication of the displayed equations by units.  The chosen
$q_j$ cut out the finite reduced polar set $P_q$ on $X$, and the slice $h$ meets
the corresponding cone line at exactly one point.  Hence this local ring is
zero-dimensional.  Since the schemes are of finite type over $\CC$, it is
Artinian and has finite length.
\end{proof}

\begin{remark}[No new degeneracy from coincident kernels]
\label{rem:nonewdegeneracy}
The coincidence of the left and right kernels in the symmetric model removes a
projective variable; it does not create an additional fiber.  Near a genuine
point the kernel line is unique and is given scheme-theoretically by the graph
$u=(-B^{-1}c,1)$.  Moreover
$(u^{\mathsf T}A_1u,\ldots,u^{\mathsf T}A_Nu)$ is a unit multiple of
$\nabla f$ on this graph.  Since $X$ is smooth, $\nabla f$ is not zero at a
genuine point.  Therefore the quadratic conormal forms cannot vanish
identically on a positive-dimensional branch through such a point.  Spurious
components elsewhere, including components over rank-drop loci or over $z=0$,
do not affect the local isolated contribution at the genuine polar points.
\end{remark}

\section{The multihomogeneous count}
\label{sec:bezout}

Let $H$ denote the hyperplane class on $\PP^N_{[z:x]}$ and $U$ the hyperplane
class on $\PP^{m-1}_{[u]}$.  The ambient product has dimension
\[
  N+(m-1)=N+m-1.
\]
The incidence system \eqref{eq:incidence} has
\[
  m+(N-2)+1=N+m-1
\]
equations, so it is square.

The $m$ equations $\Mcal u=0$ are linear in $[z:x]$ and linear in $[u]$, hence
have class $H+U$.  For a linear form
$q(Y)=\sum_i\lambda_iY_i$ on the dual coordinates,
\[
  q(u^{\mathsf T}A_1u,\ldots,u^{\mathsf T}A_Nu)
  =u^{\mathsf T}\left(\sum_i\lambda_iA_i\right)u.
\]
This is homogeneous of degree $2$ in $u$ and independent of $[z:x]$.  It is a
genuine section of $\OO(0,2)$ for every nonzero $q$ in the polar linear system.
Indeed, if $u^{\mathsf T}(\sum_i\lambda_iA_i)u$ vanished identically, then,
because the matrix $\sum_i\lambda_iA_i$ is symmetric and the characteristic is
zero, one would have $\sum_i\lambda_iA_i=0$.  Then
$\sum_i\lambda_i\partial_i f=0$ identically by differentiating the determinant.
After a linear change of the $x$-coordinates, $f$ would be independent of one
variable; since $f$ is homogeneous of degree at least two, the corresponding
coordinate point would be singular on $V(f)$, contradicting smoothness.  Hence
the polar equations have class $2U$.

There is no missing analogue of the non-symmetric redundant-left-equation
lemma.  In the ordinary determinant model one has both $\Mcal v=0$ and
$u^{\mathsf T}\Mcal=0$; one left equation is locally redundant after imposing
the right-kernel equations.  In the symmetric model there is only the single
kernel condition $\Mcal u=0$.  Locally, $m-1$ of these equations solve for the
affine kernel coordinates $u'$ and the remaining equation is the Schur
complement $g=0$.  Thus the factor is $(H+U)^m$, not $(H+U)^{m-1}$ and not a
product involving a second kernel variable.

The slice $h(z,x)=0$ has class $H$.  Therefore multihomogeneous Bezout gives
that the sum of isolated local solution multiplicities is at most
\begin{equation}
\label{eq:bezoutproduct}
  [H^N U^{m-1}]\,H(H+U)^m(2U)^{N-2}.
\end{equation}
This upper bound remains valid even if the full incidence scheme has positive-
dimensional excess components; isolated local intersection multiplicities are
bounded by the corresponding Chow-ring coefficient.

By Lemma~\ref{lem:isolated}, each of the $\pdeg_{\mathrm{top}}(X)$ genuine polar
points contributes an isolated incidence solution.  Therefore
\[
  \pdeg_{\mathrm{top}}(X)
  \le [H^N U^{m-1}]\,H(H+U)^m(2U)^{N-2}.
\]
It remains only to extract the coefficient:
\[
  (2U)^{N-2}=2^{N-2}U^{N-2}.
\]
Thus one needs the coefficient of $H^{N-1}U^{m-N+1}$ in $(H+U)^m$, namely
$\binom{m}{N-1}$.  Hence
\[
  [H^N U^{m-1}]\,H(H+U)^m(2U)^{N-2}
  =2^{N-2}\binom{m}{N-1}.
\]
This proves Theorem~\ref{thm:main}.


\section{Diagonal power sums}
\label{sec:fermat}

We first prove the arbitrary-degree diagonal corollary.  Let
\[
  F_{N,d}=x_1^d+\cdots+x_N^d,
  \qquad N\ge3,
  \qquad d\ge2.
\]
The hypersurface $X_{N,d}=V(F_{N,d})\subset\PP^{N-1}$ is smooth over $\CC$:
the partial derivatives are $d x_i^{d-1}$, and they have no common projective
zero.  Therefore
\[
  \pdeg_{\mathrm{top}}(X_{N,d})=d(d-1)^{N-2}.
\]
If $m=\sdc(F_{N,d})$, Theorem~\ref{thm:main} gives
\begin{equation}
\label{eq:diagonalineq}
  d(d-1)^{N-2}\le 2^{N-2}\binom{m}{N-1}.
\end{equation}
This is the first displayed assertion of Corollary~\ref{cor:diagonal-general}.
Using $\binom{m}{N-1}\le m^{N-1}/(N-1)!$, we obtain
\[
  m
  \ge
  \left(
  \frac{d(d-1)^{N-2}(N-1)!}{2^{N-2}}
  \right)^{1/(N-1)}.
\]
For the asymptotic form, Stirling's formula gives
\[
  ((N-1)!)^{1/(N-1)}=(1+o_N(1))\frac{N}{e}.
\]
The remaining degree factor satisfies, uniformly for $d\ge2$,
\[
  \big(d(d-1)^{N-2}\big)^{1/(N-1)}
  =(d-1)\left(\frac{d}{d-1}\right)^{1/(N-1)}
  =(1+o_N(1))(d-1),
\]
because $1\le(d/(d-1))^{1/(N-1)}\le2^{1/(N-1)}$.  Finally,
\[
  2^{(N-2)/(N-1)}=2+o_N(1).
\]
Thus
\[
  \sdc(F_{N,d})
  \ge \left(\frac{1}{2e}-o_N(1)\right)N(d-1),
\]
which proves Corollary~\ref{cor:diagonal-general}.

Taking $N=d=n$ and writing $F_n=F_{n,n}$ gives
\[
  n(n-1)^{n-2}\le 2^{n-2}\binom{\sdc(F_n)}{n-1},
\]
and hence
\[
  \sdc(F_n)\ge \left(\frac{1}{2e}-o(1)\right)n^2.
\]
This proves Corollary~\ref{cor:fermat}.

\subsection{Non-vacuity: an explicit symmetric representation}
\label{subsec:upper}

The lower bound is not merely a statement about an empty class of
representations.  We record a simple explicit construction showing
$\sdc(F_{N,d})=O(Nd)$ over $\CC$.

\begin{proposition}
\label{prop:upper}
For every $N,d\ge1$,
\[
  \sdc(F_{N,d})\le 2N(d+1)+1.
\]
In particular, for $F_n=F_{n,n}$,
\[
  \sdc(F_n)\le 2n^2+2n+1.
\]
\end{proposition}

\begin{proof}
Let $r=d+1$, and let $J$ be the $r\times r$ nilpotent matrix with ones on the
superdiagonal.  For each $i=1,\ldots,N$, set
\[
  C_i=I_r-x_iJ.
\]
Then $\det C_i=1$ and
\[
  C_i^{-1}=I_r+x_iJ+x_i^2J^2+\cdots+x_i^dJ^d.
\]
Thus
\[
  e_1^{\mathsf T}C_i^{-1}e_r=x_i^d.
\]
Define the symmetric $2r\times2r$ matrix
\[
  D_i=\begin{pmatrix}0&C_i^{\mathsf T}\\ C_i&0\end{pmatrix}.
\]
Then $\det D_i=(-1)^r$ and
\[
  D_i^{-1}=\begin{pmatrix}0&C_i^{-1}\\ (C_i^{\mathsf T})^{-1}&0\end{pmatrix}.
\]
With
\[
  b_i=\begin{pmatrix}e_1\\ e_r\end{pmatrix},
\]
one has
\[
  b_i^{\mathsf T}D_i^{-1}b_i
  =2e_1^{\mathsf T}C_i^{-1}e_r
  =2x_i^d.
\]
Let
\[
  D=\operatorname{diag}(D_1,\ldots,D_N),
  \qquad
  b=(b_1,\ldots,b_N)^{\mathsf T}.
\]
Then $D$ is symmetric affine-linear of size $2N(d+1)$,
$\det D=c\in\{\pm1\}$, and
\[
  b^{\mathsf T}D^{-1}b=2F_{N,d}.
\]
Choose $\lambda\in\CC$ satisfying $-2c\lambda^2=1$ and set
\[
  S=\begin{pmatrix}0&\lambda b^{\mathsf T}\\ \lambda b&D\end{pmatrix}.
\]
This is symmetric and affine-linear, of size $2N(d+1)+1$.  By the Schur
complement formula,
\[
  \det S
  =\det D\,\big(0-\lambda^2 b^{\mathsf T}D^{-1}b\big)
  =c(-2\lambda^2F_{N,d})=F_{N,d}.
\]
For $N=d=n$ this gives the displayed size $2n^2+2n+1$.
\end{proof}

\begin{remark}[Small-size verification of the explicit construction]
\label{rem:smallchecks}
The construction in Proposition~\ref{prop:upper} was independently expanded for
$n=2$ and $n=3$ in the specialization $F_n=F_{n,n}$.  In both cases one obtains
$b^{\mathsf T}D^{-1}b=2F_n$ and $\det S=F_n$ exactly, with sizes $13$ and $25$,
respectively.  This check is not used in the proof, but it verifies the concrete
non-vacuity construction in the first two nontrivial cases a reader is likely to
test.
\end{remark}

\section{Scope and limitations}
\label{sec:scope}

\begin{remark}[Positive characteristic]
\label{rem:positivechar}
The theorem is stated over $\CC$ for a reason.  It is false as a uniform theorem
over all fields of characteristic different from two.  If $\operatorname{char} K
=p>2$ and $n=p^r$, then the Frobenius identity gives
\[
  \sum_{i=1}^n x_i^n=(x_1+\cdots+x_n)^n.
\]
Hence
\[
  F_n=\Det\big((x_1+\cdots+x_n)I_n\big),
\]
so $\sdc(F_n)\le n$.  Also the hypersurface $V(F_n)$ is not smooth in that case.
Thus the smooth characteristic-zero hypothesis is essential for the stated
asymptotic lower bound.
\end{remark}

\begin{remark}[Exact, not border]
The argument proves an exact symmetric determinantal-complexity lower bound.  It
uses isolated solutions in a fixed incidence scheme associated with an exact
representation.  It does not prove a border symmetric determinantal-complexity
lower bound, because polar degree is not a closed invariant in arbitrary
degenerating families and the isolated local contributions can degenerate into
excess components.
\end{remark}

\begin{remark}[Relation to the non-symmetric companion preprint]
The companion preprint~\cite{SheshadriArxiv} proves, for ordinary
determinantal complexity,
\[
  \pdeg_{\mathrm{top}}(X)\le
  \sum_{a=0}^{N-2}\binom{N-2}{a}\binom{m}{N-1-a}\binom{m-1}{a}.
\]
For $F_n$, the leading coefficient of this ordinary determinant bound yields the
constant $1/(4e)$.  The symmetric incidence has only one kernel variable, so the
polar equations have class $2U$ rather than $U+V$, and the coefficient becomes
$2^{N-2}\binom{m}{N-1}$.  This is exactly the source of the constant $1/(2e)$.
\end{remark}

\section{The human--AI methodology}
\label{sec:methodology}

We describe the process honestly, both because it produced the result and
because it is itself of methodological interest.

\paragraph{Roles and models.}
A human author acted as orchestrator and domain arbiter.  Large language models
were used in separated roles: a generator proposing variants of the
polar-degree method, an adversarial critic attempting to falsify the proposed
symmetric specialization, and a verifier checking the local algebra,
multidegrees, coefficient extraction, asymptotic constant, and non-vacuity.
Convergence of independent model outputs was treated only as a signal to attempt
a human-checkable proof, not as evidence.

\paragraph{Trajectory.}
The starting point was the non-symmetric polar-degree preprint~\cite{SheshadriArxiv}.
A speculative exploration suggested that the symmetric determinant model should
have a sharper incidence count because, at a rank-$(m-1)$ symmetric matrix, the
left and right kernel lines coincide.  The present proof, however, is written so
that the symmetric local algebra and the resulting Bezout bound can be checked
without importing any lemma from that preprint.  A subsequent adversarial prompt forced the symmetric isolatedness step to be
checked from scratch: the possible failure modes were rank-drop escape,
positive-dimensional spurious components through a genuine polar point,
nonreduced scheme structure, incorrect multidegree $2U$, an erroneous Bezout
coefficient, and vacuity of the model.  The resolution is the symmetric local
normal form of Lemma~\ref{lem:isolated}; it proves that the incidence is locally
a graph and that the lifted quadratic conormal equations are the ordinary polar
equations up to units.

\paragraph{Reproducibility.}
Appendix~\ref{app:prompt} records the load-bearing verification prompt in
ASCII-normalized form.  The proof in Sections~\ref{sec:prelim}--\ref{sec:fermat}
is intended to be checked independently of its provenance.  No symbolic
computation is used as an assumption in the proof.  The direct $n=2$ and $n=3$
checks of the explicit non-vacuity construction are recorded as sanity checks in
Remark~\ref{rem:smallchecks}.

\section*{Acknowledgement of AI assistance}

The mathematical content of this paper was generated, challenged, and rewritten
through large language models operated by the human author in the adversarial
multi-model protocol described above.  The author selected the problem,
directed the verification prompts, chose the final framing, and takes
responsibility for the text.  The proofs are included in full so that the
result can be judged by ordinary mathematical standards without reference to the
models that helped produce it.

\begin{thebibliography}{99}

\bibitem{AlperBogartVelasco17}
J.~Alper, T.~Bogart, and M.~Velasco.
\newblock A lower bound for the determinantal complexity of a hypersurface.
\newblock \emph{Found. Comput. Math.} 17(3):829--836, 2017.

\bibitem{CaiChenLi10}
J.-Y.~Cai, X.~Chen, and D.~Li.
\newblock Quadratic lower bound for permanent vs. determinant in any
characteristic.
\newblock \emph{Comput. Complexity} 19(1):37--56, 2010.

\bibitem{Fulton}
W.~Fulton.
\newblock \emph{Intersection Theory}, 2nd ed.
\newblock Springer, 1998.

\bibitem{GH}
P.~Griffiths and J.~Harris.
\newblock Algebraic geometry and local differential geometry.
\newblock \emph{Ann. Sci. \'Ec. Norm. Sup\'er.} 12(3):355--452, 1979.

\bibitem{KumarVolk21}
M.~Kumar and B.~L.~Volk.
\newblock A lower bound on determinantal complexity.
\newblock In \emph{36th Computational Complexity Conference (CCC 2021)},
LIPIcs 200, article 4.  Full version: \emph{arXiv:2009.02452}.

\bibitem{MignonRessayre04}
T.~Mignon and N.~Ressayre.
\newblock A quadratic bound for the determinant and permanent problem.
\newblock \emph{Int. Math. Res. Not.} 2004(79):4241--4253, 2004.

\bibitem{Piene78}
R.~Piene.
\newblock Polar classes of singular varieties.
\newblock \emph{Ann. Sci. \'Ec. Norm. Sup\'er.} 11(2):247--276, 1978.

\bibitem{SheshadriArxiv}
K.~Sheshadri.
\newblock A near-quadratic lower bound on the determinantal complexity of
$\sum_i x_i^n$ via the polar degree of the tangent cone.
\newblock arXiv:7680505, 2026.

\bibitem{Valiant79}
L.~G.~Valiant.
\newblock Completeness classes in algebra.
\newblock In \emph{Proc. 11th ACM Symposium on Theory of Computing},
pp.~249--261, 1979.

\bibitem{Yabe15}
A.~Yabe.
\newblock Bi-polynomial rank and determinantal complexity.
\newblock \emph{arXiv:1504.00151}, 2015.

\end{thebibliography}

\appendix

\section{Verification prompt}
\label{app:prompt}

The following prompt was used to force a referee-level settlement of the
symmetric specialization.  It is reproduced in ASCII-normalized form and lightly
line-wrapped; mathematical symbols such as ``not equal'' and ``Omega'' were
normalized to avoid LaTeX encoding issues.

\begin{verbatim}
Establish or refute, to journal-referee rigor, the following claimed theorem in
algebraic complexity, working it out fully yourself rather than deferring to any
prior derivation.

CLAIM. Define symmetric determinantal complexity sdc(f) as the least m such that
f = det(M(x)) for a SYMMETRIC m x m matrix M of affine-linear forms. Then for
F_n = sum_i x_i^n,
    sdc(F_n) >= (1/2e - o(1)) n^2,
via the bound delta_top(X) <= 2^{N-2} C(m, N-1) for any size-m symmetric
determinantal representation of a smooth degree-d hypersurface X = V(f) in
P^{N-1}, where delta_top = d(d-1)^{N-2}.

The intended mechanism: at a smooth point the symmetric M has rank m-1 with
COINCIDING left/right kernel, so the incidence lives in P^N_[z:x] x P^{m-1}_[u]
with one kernel variable; the conormal identity is [partial_i f] =
[u^T A_i u]; the polar equations q_j(u^T A_i u)=0 have class 2U;
multihomogeneous Bezout then gives
    2^{N-2} C(m, N-1) = [H^N U^{m-1}] H (H+U)^m (2U)^{N-2}.

Do NOT assume the count is valid. The entire claim rests on a single
load-bearing step, and your job is to settle that step, not to restate the
mechanism. Treat a confirmation and a refutation with equal suspicion; prior
versions of the analogous (non-symmetric) step were internally convincing and
WRONG.

Settle each of the following with a definite verdict -- HOLDS (full proof),
FAILS (explicit verified counterexample), or UNRESOLVED (state exactly what is
missing). Do not round UNRESOLVED toward either side.

1. SYMMETRIC ISOLATEDNESS. Does each genuine polar point lift to a
   ZERO-DIMENSIONAL, scheme-theoretically isolated solution of the symmetric
   incidence system {M(x)u=0, q_j(u^T A_i u)=0, slice}? Prove it via the
   symmetric local normal form: with M = [[B,c],[c^T,s]], det B a unit near the
   point, kernel u=(-B^{-1}c, 1), and det M = (det B)(s - c^T B^{-1} c).
   Verify that the Schur complement identification goes through WITH THE
   SYMMETRY CONSTRAINT -- i.e. that restricting to symmetric M does not destroy
   the graph structure, the unit scaling adj M = alpha u u^T, or the
   elimination to the polar slice. State explicitly whether coincidence of the
   two kernels introduces any NEW degeneracy (e.g. u^T A_i u vanishing
   identically on a positive-dimensional locus) absent in the non-symmetric
   case.

2. THE MULTIDEGREE 2U. Confirm q_j(u^T A_i u) is genuinely class 2U and that
   the symmetric incidence is a SQUARE system on P^N x P^{m-1} (count equations
   vs dimensions). Verify no analogue of the redundant-left-equation reduction
   is silently needed or silently missing.

3. THE BEZOUT EXTRACTION. Verify
   [H^N U^{m-1}] H (H+U)^m (2U)^{N-2} = 2^{N-2} C(m, N-1) exactly, and that the
   root extraction yields the constant 1/(2e), not 1/(4e) or something else.

4. EXISTENCE / NONDEGENERACY. Does a smooth symmetric determinantal
   representation of F_n even exist for the relevant m, and is the generic
   polar section reduced of the right cardinality in the symmetric setting? If
   symmetric representations of F_n are obstructed, the theorem is vacuous or
   false -- check this.

Close with: (i) overall verdict on the n^2 bound and the 1/(2e) constant; (ii)
the single weakest point; (iii) an explicit list of anything asserted but not
fully verified.
\end{verbatim}

\section{Checklist of load-bearing verifications}
\label{app:checklist}

For convenience, we spell out where each possible failure mode is addressed.

\begin{enumerate}[leftmargin=2em]
\item Rank-drop escape cannot pass through a genuine point because the local
normal form is taken on the open set $z\ne0$ and $\det B\ne0$, and the local
incidence is a graph over $g=0$; see Lemma~\ref{lem:isolated}.
\item A positive-dimensional spurious component cannot pass through a genuine
point because on the local graph the lifted polar equations are ordinary polar
equations multiplied by units; see \eqref{eq:unitpolar}.
\item Nonreduced structure is controlled because the completed local incidence
ring is Artinian.  Nonreduced isolated multiplicity is allowed and is exactly
what Bezout counts.
\item The polar equations have class $2U$ because $u^{\mathsf T}A_i u$ is a
quadratic form in the single projective kernel variable.  No left-kernel
reduction is present in the symmetric system.
\item The coefficient extraction is the elementary identity
\[
  [H^N U^{m-1}]H(H+U)^m(2U)^{N-2}=2^{N-2}\binom{m}{N-1}.
\]
\item The model is nonempty for $F_n$ by Proposition~\ref{prop:upper}.
\end{enumerate}

\end{document}
